\documentclass[11pt]{article}

\usepackage[T1]{fontenc}
\usepackage[utf8]{inputenc}
\usepackage{lmodern}
\usepackage[margin=1in]{geometry}
\usepackage{microtype}
\usepackage{amsmath,amssymb,mathtools,bm}
\usepackage{booktabs,tabularx,array}
\usepackage{enumitem}
\usepackage{xcolor}
\usepackage{hyperref}

\hypersetup{
  colorlinks=true,
  linkcolor=blue!55!black,
  urlcolor=blue!65!black,
  citecolor=blue!55!black,
  pdftitle={BCC Binary-Junction Node Splitting in ParaDiS}
}

\setlist{nosep,leftmargin=*}
\newcommand{\code}[1]{\texttt{\detokenize{#1}}}
\newcommand{\vect}[1]{\bm{#1}}

\title{BCC Binary-Junction Node Splitting in ParaDiS}
\author{ParaDiS-Extended Development Notes}
\date{August 2026}

\begin{document}

\maketitle

\begin{abstract}
This note documents the implementation of three-arm BCC binary-junction
splitting in the serial and parallel ParaDiS topology paths.  The port follows
the ExaDiS/OpenDiS implementation pinned at commit
\code{351607fa0f04aa34d29ee253438edde994473fd6}, with adaptations for ParaDiS
node constraints, remeshing, and its serialized cross-domain operation list.
It also describes the supplied one-step regression fixture and the topology
expected after a successful split.
\end{abstract}

\tableofcontents

\section{Scope and reference implementation}

The change extends the existing multi-node separation machinery, which
previously considered nodes with at least four arms, to one physically
specific degree-three case: a non-planar BCC binary junction.  Both topology
selectors use the same classifier, arm-set restriction, mobility test,
acceptance threshold, split geometry, and corner-node semantics.

The source behavior was compared against the pinned upstream topology code:

\begin{center}
\url{https://github.com/llnl/exadis/tree/351607fa0f04aa34d29ee253438edde994473fd6/src/topology_types}
\end{center}

The algorithm is ported rather than copied mechanically.  ParaDiS has a
different ownership model and communicates topology changes through a packed
remote-operation stream, so explicit propagation of the new node constraint
is required here.

\section{Binary-junction recognition}

\subsection{Topological and Burgers-vector tests}

The shared function
\code{BCC_binary_junction_node(Home_t*, Node_t*, real8[3], int*)} rejects a
node unless it has exactly three arms.  For arm $i$, the minimum-image line
direction is
\begin{equation}
  \vect{t}_i =
  \frac{\operatorname{ZImage}(\vect{x}_i-\vect{x}_0)}
       {\left\lVert\operatorname{ZImage}(\vect{x}_i-\vect{x}_0)\right\rVert}.
\end{equation}
Degenerate arms and missing neighbors cannot form a valid candidate.  An arm
is classified as the junction arm when
\begin{equation}
  \lVert\vect{b}_i\rVert > 1 + 10^{-12}.
\end{equation}
The node is a binary junction only when exactly one arm passes this test and
the remaining two are non-junction arms (called glide arms by the topology
implementation).  This magnitude test does not independently validate a
glissile Burgers-vector family.  The classifier returns the junction arm
index and its unit line direction $\vect{t}_{\mathrm{j}}$.

Candidate construction adds the following guards:
\begin{itemize}
  \item the material must be BCC and \code{split3node} must be enabled;
  \item planar junctions are excluded;
  \item surface and pinned nodes are excluded; and
  \item the existing short-segment filter must pass.
\end{itemize}

\subsection{Five-degree planarity test}

For each glide arm, the parent-plane normal is formed as
\begin{equation}
  \vect{n}_i = \vect{b}_i \times \vect{t}_i.
\end{equation}
An arm is treated as edge-bearing when
$\vect{n}_i\mathbin{\cdot}\vect{n}_i>10^{-2}$.  With
$\theta_{\mathrm{tol}}=5^{\circ}$, the junction is classified as planar as
follows:

\begin{itemize}
  \item For two edge-bearing parents, normalize
        $\vect{d}=\vect{n}_1\times\vect{n}_2$ and require
        $\bigl|\,|\vect{d}\cdot\vect{t}_{\mathrm{j}}|-1\,\bigr|
        <1-\cos\theta_{\mathrm{tol}}$.
  \item For one edge-bearing parent, require
        $|\vect{n}_1\cdot\vect{t}_{\mathrm{j}}|<
        \sin\theta_{\mathrm{tol}}$.
  \item If both parents are screw-like, the junction is planar by
        convention.
\end{itemize}

Only a recognized junction for which this planarity result is false proceeds
to split evaluation.

\section{Split construction and acceptance}

\subsection{Allowed arm set}

\code{SMNGetArmSets()} now supports degree-three nodes.  It produces the three
possible two-arm selections, while preserving the normal combinatorial
enumeration for larger nodes.  For a binary junction, evaluation retains only
the selection containing the two glide arms; the junction arm must remain on
the other trial node.  The split therefore produces one degree-three node and
one degree-two node connected by the newly created segment.

\subsection{Trial mobility and geometry}

The trial split initially creates two coincident nodes and reuses the saved
arm forces.  Mobility is evaluated for both nodes, the existing junction
force/velocity adjustment is applied, and the velocity of the degree-two
trial node is set to zero.  This is the binary-junction boundary condition
used by the reference implementation.  Mobility return codes are combined so
an error from either trial node rejects the candidate.

The candidate nodes are then separated along the accepted trial direction.
The required total separation is
\begin{equation}
  d_{\mathrm{split}} = 2r_{\mathrm{ann}} + 10^{-12}.
\end{equation}
If both nodes move, each moves by half this distance; otherwise the mobile
node moves by the full distance.  Normalization in the opposite-velocity
case is guarded against a zero norm, and the existing glide-plane rule is
retained.  The evaluated trial positions are saved and reused for the final
three-arm split instead of being recomputed by the general
\code{CalcJunctionLen()} path.  The small positive offset prevents the new
connector from lying exactly on the immediate remesh threshold.

\subsection{Energetic and kinematic acceptance}

Let $\vect{f}_k$ and $\vect{v}_k$ be the force and velocity of trial node
$k$.  The split power is evaluated as
\begin{equation}
  P_{\mathrm{s}} =
  \vect{f}_1\!\cdot\!\vect{v}_1+
  \vect{f}_2\!\cdot\!\vect{v}_2-
  v_{\mathrm{noise}}
  \left(\lVert\vect{f}_1\rVert+\lVert\vect{f}_2\rVert\right),
\end{equation}
where
\begin{equation}
  v_{\mathrm{noise}} =
  \code{splitMultiNodeAlpha}\,
  \frac{\code{rTol}}{\code{deltaTT}}.
\end{equation}
The relative velocity must point along the proposed separation direction.
The binary-junction candidate is accepted only if it improves upon the best
power by more than the numerical tolerance and also satisfies
$P_{\mathrm{s}}>1$.  The latter threshold is specific to degree-three
candidates and prevents marginal splits from being selected.

\section{Serial topology path}

The serial implementation resides primarily in \path{src/Topology.cc}:

\begin{enumerate}
  \item \code{SplitMultiNodes()} admits degree-three candidates only through
        the shared BCC classifier and the eligibility guards above.
  \item The existing node/neighbor backup and temporary \code{SplitNode()}
        workflow evaluates the single meaningful glide-arm separation.
  \item The degree-two trial velocity, split-power threshold, and saved trial
        positions apply only to the binary-junction case; behavior for nodes
        with four or more arms remains on the established path.
  \item After the final split succeeds, whichever result has degree two is
        marked \code{CORNER_NODE} and its complete constraint value is queued
        for remote domains when the operation is global.
\end{enumerate}

Temporary arm storage was enlarged to hold both selected glide arms for a
three-arm split, and temporary reposition flags are initialized before each
evaluation.  These details avoid under-allocation and stale-state behavior
when the new, smaller topology enters code originally sized for nodes of
degree four or greater.

\section{Parallel topology path}

The parallel selector continues to use the staged
\code{SplitMultiNodesParallel()} workflow:

\begin{enumerate}
  \item \code{BuildGhostMultiNodeList()} and
        \code{BuildLocalMultiNodeList()} call the same splittability predicate
        for native and ghost degree-three nodes.
  \item \code{SMNEvalSplitForces()} constructs the eligible trial and computes
        each domain's partial segment forces.
  \item The existing force-buffer exchange reduces those contributions onto
        the owning domain.
  \item \code{SMNEvalMultiNodeSplits()} performs the final mobility,
        separation, power, and direction tests and executes the winning split.
\end{enumerate}

All three degree-three split-info slots remain present in the communication
layout, but the two arm sets containing the junction arm are skipped.  Keeping
the full slot indexing makes force packing and unpacking consistent across
domains.  The parallel evaluation mirrors the serial fixes: adequate arm and
force storage, combined mobility status, guarded velocity normalization,
zero degree-two velocity, the $P_{\mathrm{s}}>1$ threshold, saved final
positions, and corner marking after a successful split.

In a non-\code{PARALLEL} build, initialization forces
\code{useParallelSplitMultiNode} to zero.  Thus the parallel selector is only
meaningful when ParaDiS was compiled with parallel topology support.

\section{Corner preservation and remeshing}

\code{CORNER_NODE} is the constraint bit \code{0x00000010} (decimal 16).  It
identifies the physical degree-two corner created by a successful binary
junction split.  The bit is included in the valid constraint mask.  During
topology initialization and native/ghost candidate-list construction, stale
corner markers are removed from nodes observed with a degree other than two.

Both remesh coarsening paths apply the same length rule.  An edge touching a
corner may be merged only when
\begin{equation}
  L \leq 2r_{\mathrm{ann}}.
\end{equation}
The candidate phase avoids selecting a corner whose two attached edges are
both longer than that limit.  The execution phase gates the two neighboring
edges independently, so arbitrary arm ordering cannot cause a long
corner-attached edge to collapse merely because the other edge is short.
Very short pathological segments may still be removed.

\section{Cross-domain constraint propagation}

ParaDiS topology operations are replayed on neighboring MPI domains before
the immediate remesh stage.  The new append-only operation
\code{REM_OP_RESET_CONSTRAINT} carries a node tag and the complete integer
constraint value:

\begin{enumerate}
  \item \code{SplitNode()} first emits its existing create/connectivity
        operations.
  \item The successful serial or parallel caller marks the degree-two node
        and appends \code{AddOpResetConstraint()}.
  \item \code{FixRemesh()} replays the byte stream in order, so a newly created
        ghost exists before its constraint is reset.
\end{enumerate}

Appending the enumeration value preserves the numeric identifiers of all
pre-existing remote operations.  Nevertheless, the packed operation stream
is an executable-level protocol: every MPI rank must use the same rebuilt
binary.

\section{Control parameters}

\begin{table}[ht]
\centering
\begin{tabularx}{\textwidth}{@{}>{\ttfamily}l X@{}}
\toprule
\normalfont Parameter & Meaning \\
\midrule
split3node & Boolean switch for eligible BCC three-arm binary-junction
splitting.  The default is 1 (enabled), and input is normalized to zero or
one. \\
useParallelSplitMultiNode & Selects the topology implementation: 0 uses
\code{SplitMultiNodes()}, while 1 uses \code{SplitMultiNodesParallel()} in a
\code{PARALLEL} build. \\
splitMultiNodeFreq & Existing cadence controlling which cycles evaluate
multi-node splits. \\
rann & Existing annihilation distance; the new split and corner-remesh
thresholds are defined from $2r_{\mathrm{ann}}$. \\
\bottomrule
\end{tabularx}
\end{table}

No new material type or mobility law is introduced.  The feature activates
only when the existing material selection is BCC.

\section{Implementation map}

\begin{table}[ht]
\centering
\small
\begin{tabularx}{\textwidth}{@{}>{\raggedright\arraybackslash}p{0.35\textwidth} X@{}}
\toprule
Files & Responsibility \\
\midrule
\path{include/Node.h}, \path{src/Initialize.cc} & Define and validate
\code{CORNER_NODE}. \\
\path{include/Param.h}, \path{src/Param.cc},
\path{src/InputSanity.cc}, \path{src/ParadisInit.cc} & Bind, normalize, and
report \code{split3node}. \\
\path{include/Topology.h}, \path{src/Topology.cc} & Shared classifier and
serial split implementation. \\
\path{src/SMNSupport.cc}, \path{src/SplitMultiNodesParallel.cc} & Parallel
candidate lists, arm sets, and orchestration. \\
\path{src/SMNEvalSplitForces.cc},
\path{src/SMNEvalMultiNodeSplits.cc} & Distributed trial-force evaluation and
owning-domain split selection. \\
\path{src/Remesh.cc} & Corner-aware candidate and per-edge coarsening gates. \\
\path{include/OpList.h}, \path{src/RemoteOps.cc}, \path{src/FixRemesh.cc} &
Serialize, emit, and replay full constraint resets. \\
\path{tests/bcc_binary_junction_node.ctrl},
\path{tests/bcc_binary_junction_node.data} & One-step regression control and
four-node input network. \\
\bottomrule
\end{tabularx}
\end{table}

\section{Regression fixture}

\subsection{Input geometry}

The data file contains four nodes in a periodic box.  The central node has
one $\langle100\rangle$ junction arm of magnitude approximately 1.1547 and
two $\tfrac12\langle111\rangle$ glide arms of unit magnitude.  Their line
directions make the junction non-planar under the five-degree classifier.
The three endpoint nodes are pinned so the one-step test isolates the
topology operation from endpoint motion.

The control file selects \code{BCC_0}, enables \code{split3node}, evaluates
multi-node splitting every cycle, applies a $zz$ stress of
$5\times10^8$, and runs one step.  It defaults to the serial selector.

\subsection{Commands}

Run from the repository's \path{tests} directory:

\begin{verbatim}
../bin/paradis -d bcc_binary_junction_node.data \
    bcc_binary_junction_node.ctrl
\end{verbatim}

Set \code{useParallelSplitMultiNode = 1} in the control file to exercise the
parallel topology selector.  For the MPI ownership/remote-operation check,
use a \code{PARALLEL} build and run:

\begin{verbatim}
mpirun -n 2 ../bin/paradis -doms 2 1 1 \
    -d bcc_binary_junction_node.data \
    bcc_binary_junction_node.ctrl
\end{verbatim}

\subsection{Expected result}

After one step, the restart under
\path{bcc_binary_junction_node_results/restart/restart.data} must contain:

\begin{itemize}
  \item five nodes in total;
  \item the three original pinned degree-one endpoints;
  \item exactly one degree-two node with constraint 16
        (\code{CORNER_NODE}); and
  \item the surviving degree-three node without the corner constraint.
\end{itemize}

The serial and parallel one-rank paths should produce the same restart
topology.  The two-rank run additionally checks that the degree-two corner
constraint survives ownership and ghost replay before remeshing.

\subsection{Validation performed}

The final regression review exercised both complete build configurations and
all relevant topology/ownership combinations:

\begin{itemize}
  \item serial and \code{PARALLEL} builds completed successfully;
  \item one-rank runs with topology selectors 0 and 1 produced byte-identical
        restart data;
  \item two-rank MPI runs completed with both topology selectors; and
  \item the final restart contained the five-node topology and unique
        degree-two constraint-16 corner described above.
\end{itemize}

\section{Compatibility boundaries}

This implementation deliberately does not generalize three-arm splitting.
It applies only to BCC nodes matching the exact binary-junction signature and
rejects planar, pinned, surface, degenerate, and short-segment candidates.
Nodes with four or more arms retain their prior multi-node split behavior.
Because \code{CORNER_NODE} is meaningful only for degree-two nodes, topology
list construction removes the bit from any node whose degree later changes.

\end{document}
